% META: {"id": "1NSI-AGT-CO-003", "chapitre": "1NSI-ALGO-PARCOURS-TRIS", "type_objet": "corrige", "exercice_ref": "1NSI-AGT-EX-003", "capacites": ["P-ALGO-02A", "P-ALGO-02C"], "sources_inspiration": [], "status": "verified"}
% BEGIN-VERIFY
% def tri_insertion(tableau):
%     tableau = tableau[:]
%     for i in range(1, len(tableau)):
%         valeur = tableau[i]
%         j = i - 1
%         while j >= 0 and tableau[j] > valeur:
%             tableau[j + 1] = tableau[j]
%             j -= 1
%         tableau[j + 1] = valeur
%     return tableau
% depart = [4, 2, 7, 1]
% assert tri_insertion(depart) == [1, 2, 4, 7]
% END-VERIFY
\begin{corrige}{1NSI-AGT-EX-003}
\textbf{1.} Tableau initial : \lstinline{[4, 2, 7, 1]}.
\begin{itemize}
  \item Après $i=1$ (insertion de $2$) : \lstinline{[2, 4, 7, 1]}
  \item Après $i=2$ (insertion de $7$, déjà à sa place) : \lstinline{[2, 4, 7, 1]}
  \item Après $i=3$ (insertion de $1$) : \lstinline{[1, 2, 4, 7]}
\end{itemize}

\textbf{2.} Après $i=1$ : \lstinline{[2, 4]} est trié. Après $i=2$ : \lstinline{[2, 4, 7]}
est trié. Après $i=3$ : \lstinline{[1, 2, 4, 7]} est trié. L'invariant est vérifié à
chaque étape.

\textbf{3.} Oui, le tableau final \lstinline{[1, 2, 4, 7]} est correctement trié, ce que
confirme l'exécution du code.
\end{corrige}
